% applications/railway.tex

\chapter{Railway Alignment Applications}

\label{chap:applications-railway}

Railway engineering is the primary application domain used in this
thesis to demonstrate AIM capabilities. The chapter applies established
alignment semantics and state/runtime concepts to practical railway
questions.

% ============================================================
\section{Application Context}
% ============================================================

Railway alignment use involves coupled horizontal geometry, vertical
profile, cant, velocity, admissibility constraints, and maintenance
context. In AIM these are handled as coordinated layers consumed from
prior parts.

% ============================================================
\section{Use of Canonical State and Runtime Concepts}
% ============================================================

Railway applications consume canonical pose2/pose3 and runtime services
rather than redefining them. Operational quantities are evaluated as
derived runtime quantities from established operators.

This supports engineering tasks such as alignment comparison,
transition-quality assessment, and realization-aware evaluation in
existing infrastructure.

% ============================================================
\section{Transition and Coupled Behaviour in Practice}
% ============================================================

Transition behaviour is interpreted through curvature evolution together
with cant and operating speed context. The relevant application-level
focus is not only geometric boundary matching, but engineering
consequence in comfort, admissibility, and lifecycle behaviour.

% ============================================================
\section{Operational and Maintenance Consequences}
% ============================================================

In application workflows, alignment decisions interact with constraints
on cant, transition quality, and maintainability. AIM supports this by
keeping persistent semantics separate from derived operational
quantities and by enabling runtime/realization-aware evaluation.

% ============================================================
\section{Interoperability Context}
% ============================================================

Railway applications often require exchange with BIM/GIS and format
workflows. In this part, such standards are treated as interoperability
mechanisms downstream of canonical engineering semantics.

% ============================================================
\section{Connection to AIM}
% ============================================================

\begin{summary}
Railway applications in AIM demonstrate how alignment semantics, state,
runtime evaluation, and realization-aware interpretation jointly support
practical engineering decisions.

The chapter remains application-focused and does not redefine canonical
concepts owned by upstream parts.
\end{summary}
